\begin{lemma}[A good hierarchy survives outside-chain recoloring]
Let two corresponding ETT presentations have the same terminal tree,
stored sets, starred core, ordered stages, and dividers.  Require also
that they have the same PE/non-PE terminal status and, in the common PE
case, the same chosen closure.  Let the first presentation be under
$\varphi_n$ and the second under $\sigma_n$.  Suppose the first
presentation is a good hierarchy and $\sigma_n$ is obtained from
$\varphi_n$ by recoloring finitely many $(\alpha,\beta)$-chains fully
outside the terminal tree.  Then the second presentation is a good
hierarchy with the same $D_{n,j}$ at corresponding stages and exactly
the same $\Gamma^j_\eta$ for every $\eta\in D_{n,j}$.  Values of the
total Gamma functions outside these paper domains are irrelevant.
\end{lemma}
\begin{proof}
By \noderef{OutsideKempeRecoloringStable}, every edge incident with the
terminal tree has the same value under the two colorings, and every
vertex of the terminal tree has the same missing-color set.  Every
hierarchy stage and the starred base are vertex subsets of the terminal
tree by \noderef{TashkinovHierarchy}, \noderef{HierarchyBase}, and the
equalities in \noderef{SameHierarchyStages}.  In the PE case the chosen
closure is also equal in the two bases and is contained in the common
starred core by \noderef{OrderedTreeUnion}; in a non-PE case it enters
none of the five good-hierarchy clauses.  Therefore all paper-relevant
missing-color sets, all colors on their displayed edge differences, and
every colored boundary-edge set agree under $\varphi_n$ and $\sigma_n$.

The raw stored sets agree at every rung by hypothesis.  Combining this
with the preceding equality of missing-color sets in the definition
\noderef{HierarchyStoredColorsAt} gives the same $D_{n,j}$ at every
stage.  Keep the original two-color sets
$\Gamma^j_\eta$ as the Gamma field for the second hierarchy.

We check the five fields of \noderef{GoodHierarchy}.  Pair cardinality
and pairwise disjointness are literal properties of the unchanged finite
sets.  Clause (i) is unchanged because both the stage-missing sets and
the segment edge-color sets of \noderef{HierarchySegmentColors} agree.
Clause (iii) is unchanged because the stage differences have the same
vertices and missing colors, and every unrefreshed Gamma equality is
literal.  Clause (iv) is unchanged because the starred base,
$T_n$, $R_n$, their common missing-color set, and all relevant boundary
colors agree.  Clause (v) is unchanged for the same reason: both the
missing-color set and every colored boundary identity of the stage agree.
Thus all five clauses hold under $\sigma_n$.

The constructed witness has the same paper-extensional hierarchy
presentation by the assumed \noderef{SameHierarchyStages} relation.
That relation includes equality of PE status and, in the PE case,
equality of the closure; in the common non-PE case it ignores both unused
closure fields.  At each stage, the equality of raw stored colors and of
stage missing-color sets identifies the corresponding $D_{n,j}$ sets.
On every color in this common set the copied Gamma value is literally the
old Gamma value.  These are exactly the transported, stagewise clauses of
\noderef{SameGoodHierarchy}; no condition is needed at a color outside
$D_{n,j}$.
\end{proof}
